%! TEX root = master.tex
\chapter{3}{Singularities, Laurent Series and Residues}

Near an isolated singularity a function is no longer a power series, but it is
still a series in both positive and negative powers. The negative part
classifies the singularity, and a single coefficient, the residue, is all
that survives when we integrate around it.

\section{Laurent series}\label{sec:laurent}

%! TEX root = ../master.tex
\begin{theorem}[Laurent's theorem]\label{thm:laurent}
  If $f$ is analytic in the annulus $r < \abs{z - z_0} < R$, then
  \[
    f(z) = \sum_{n=-\infty}^{\infty} a_n (z - z_0)^n
  \]
  there, where both halves of the series converge absolutely in the annulus
  and uniformly on every closed subannulus. The coefficients are unique, and
  \[
    a_n = \frac{1}{2\pi i} \oint_C \frac{f(\zeta)}{(\zeta - z_0)^{n+1}} \dd \zeta
  \]
  for any positively oriented circle $C \colon \abs{\zeta - z_0} = \rho$ with
  $r < \rho < R$.
  \thmsource{\cite[ch.~5]{saff-snider}}
\end{theorem}


\begin{definition}\label{def:singularities}
  Let $f$ be analytic in $0 < \abs{z - z_0} < R$, with Laurent series
  $\sum_{n=-\infty}^{\infty} a_n (z - z_0)^n$ there. The isolated singularity
  $z_0$ is
  \begin{itemize}
    \item \emph{removable} if $a_n = 0$ for every $n < 0$;
    \item a \emph{pole of order $m$} if $a_{-m} \neq 0$ and $a_n = 0$ for
      every $n < -m$;
    \item \emph{essential} if $a_n \neq 0$ for infinitely many $n < 0$.
  \end{itemize}
  The \emph{residue} of $f$ at $z_0$ is $\Res(f; z_0) = a_{-1}$.
\end{definition}

\begin{exercise}\label{ex:exp-inverse-square}
  Find the Laurent series of $f(z) = e^{1/(2z^2)}$ about $0$, classify the
  singularity, and find $\Res(f; 0)$.
\end{exercise}

\begin{solution}
  The series $e^w = \sum_{j=0}^{\infty} w^j / j!$ converges for every $w$,
  because $e^w$ is entire. Putting $w = 1/(2z^2)$ gives, for every $z \neq 0$,
  \[
    e^{1/(2z^2)}
    = \sum_{j=0}^{\infty} \frac{1}{j! \, 2^j z^{2j}}
    = 1 + \frac{1}{2z^2} + \frac{1}{8z^4} + \frac{1}{48z^6} + \dotsb.
  \]
  By uniqueness of Laurent coefficients (\cref{thm:laurent}), this is the
  Laurent series on $0 < \abs{z} < \infty$. It has infinitely many negative
  powers, so $0$ is an essential singularity. Only even powers occur, so
  $a_{-1} = 0$ and $\Res(f; 0) = 0$.
\end{solution}

\begin{remark}
  A zero residue does not mean a mild singularity: here it sits beside an
  essential singularity. It only says that $\oint f = 0$ around small circles.
\end{remark}

\section{Removable singularities}\label{sec:removable}

%! TEX root = ../master.tex
\begin{theorem}[Removable singularities]\label{thm:removable}
  Let $f$ have an isolated singularity at $z_0$. The following are equivalent:
  \begin{enumerate}[label=(\roman*)]
    \item $z_0$ is a removable singularity;
    \item $\lim_{z \to z_0} f(z)$ exists and is finite;
    \item $f$ is bounded in some punctured disc $0 < \abs{z - z_0} < \delta$.
  \end{enumerate}
  In that case, setting $f(z_0) = \lim_{z \to z_0} f(z)$ makes $f$ analytic at
  $z_0$.
  \thmsource{\cite[ch.~5]{saff-snider}}
\end{theorem}


\begin{exercise}\label{ex:one-minus-cos}
  Show that $f(z) = \dfrac{1 - \cos z}{z^2}$ has a removable singularity at
  $0$, and find the value of $f(0)$ that makes $f$ analytic there.
\end{exercise}

\begin{solution}
  For every $z$,
  $\cos z = 1 - \frac{z^2}{2!} + \frac{z^4}{4!} - \frac{z^6}{6!} + \dotsb$,
  so for $z \neq 0$
  \[
    f(z) = \frac{1}{2!} - \frac{z^2}{4!} + \frac{z^4}{6!} - \dotsb
    = \frac12 - \frac{z^2}{24} + \frac{z^4}{720} - \dotsb.
  \]
  The Laurent series has no negative powers, so the singularity is removable.
  Equivalently, $\lim_{z \to 0} f(z) = \frac12$ (\cref{thm:removable}).
  Setting $f(0) = \frac12$ makes $f$ the sum of an everywhere convergent power
  series, so $f$ is entire.
\end{solution}

\section{Zeros and poles}\label{sec:poles}

%! TEX root = ../master.tex
\begin{theorem}[Zeros and poles]\label{thm:poles}
  Let $m \geq 1$.
  \begin{enumerate}[label=(\roman*)]
    \item $f$ has a pole of order $m$ at $z_0$ if and only if
      $f(z) = g(z) / (z - z_0)^m$ near $z_0$, with $g$ analytic at $z_0$ and
      $g(z_0) \neq 0$.
    \item $g$ has a zero of order $m$ at $z_0$ if and only if
      $g(z) = (z - z_0)^m h(z)$ with $h$ analytic at $z_0$ and $h(z_0) \neq 0$.
    \item If $g$ has a zero of order $m$ at $z_0$, then $1/g$ has a pole of
      order $m$ there.
  \end{enumerate}
  \thmsource{\cite[ch.~5]{saff-snider}}
\end{theorem}


\begin{exercise}\label{ex:double-poles}
  Find the poles of $f(z) = \dfrac{e^z}{z^2 (z - 2)^2}$ and their orders.
\end{exercise}

\begin{solution}
  $e^z$ is entire, so $f$ is analytic except at the zeros $0$ and $2$ of the
  denominator.

  \emph{At $0$.} $f(z) = g(z) / z^2$ with $g(z) = e^z / (z - 2)^2$, which is
  analytic at $0$ with $g(0) = \frac14 \neq 0$. By \cref{thm:poles}, $0$ is a
  pole of order $2$.

  \emph{At $2$.} $f(z) = g(z) / (z - 2)^2$ with $g(z) = e^z / z^2$, analytic
  at $2$ with $g(2) = e^2 / 4 \neq 0$. So $2$ is also a pole of order $2$.
\end{solution}

\section{Residues}\label{sec:residues}

%! TEX root = ../master.tex
\begin{theorem}[Cauchy's residue theorem]\label{thm:residue-theorem}
  Let $\Gamma$ be a simple closed positively oriented contour, and let $f$ be
  analytic inside and on $\Gamma$ except at finitely many points
  $z_1, \dots, z_n$ inside $\Gamma$. Then
  \[
    \oint_\Gamma f(z) \dd z = 2\pi i \sum_{k=1}^{n} \Res(f; z_k).
  \]
  \thmsource{\cite[ch.~6]{saff-snider}}
\end{theorem}

\begin{proof}
  Surround each $z_k$ by a small positively oriented circle $C_k$, with the
  discs disjoint and inside $\Gamma$. By \cref{thm:deformation},
  $\oint_\Gamma f = \sum_k \oint_{C_k} f$. On $C_k$, $f$ equals its Laurent
  series about $z_k$, which converges uniformly there (\cref{thm:laurent}),
  so it may be integrated term by term. Since
  \[
    \oint_{C_k} (z - z_k)^n \dd z =
    \begin{cases} 2\pi i, & n = -1, \\ 0, & n \neq -1, \end{cases}
  \]
  only the $a_{-1}$ term survives, and $\oint_{C_k} f = 2\pi i \Res(f; z_k)$.
\end{proof}

%! TEX root = ../master.tex
\begin{theorem}[Residue at a pole]\label{thm:residue-at-pole}
  If $f$ has a pole of order $m$ at $z_0$, then
  \[
    \Res(f; z_0) = \frac{1}{(m-1)!} \lim_{z \to z_0}
      \frac{\dd^{m-1}}{\dd z^{m-1}} \bigl[ (z - z_0)^m f(z) \bigr].
  \]
  For a simple pole, $\Res(f; z_0) = \lim_{z \to z_0} (z - z_0) f(z)$.
  \thmsource{\cite[ch.~6]{saff-snider}}
\end{theorem}

\begin{proof}
  Near $z_0$,
  \[
    (z - z_0)^m f(z)
    = a_{-m} + a_{-m+1}(z - z_0) + \dots + a_{-1}(z - z_0)^{m-1} + \dotsb,
  \]
  a power series once its removable singularity at $z_0$ is filled in. By
  \cref{thm:uniqueness}, its $(m-1)$st derivative at $z_0$ is
  $(m-1)! \, a_{-1}$, and this derivative is continuous, so it equals the limit.
\end{proof}


\begin{exercise}\label{ex:z-squared-sin}
  Let $f(z) = \dfrac{1}{z^2 \sin z}$. Find $\Res(f; 0)$ in two ways, first by
  \cref{thm:residue-at-pole} and then from the Laurent series, and use it to
  evaluate $\oint_{\abs{z} = 1} f(z) \dd z$.
\end{exercise}

\begin{solution}
  \emph{Order of the pole.} From the series for $\sin z$,
  \[
    z^2 \sin z = z^3 h(z),
    \qquad h(z) = 1 - \frac{z^2}{6} + \frac{z^4}{120} - \dotsb,
  \]
  with $h(0) = 1 \neq 0$. So $z^2 \sin z$ has a zero of order $3$ at $0$, and
  $f$ has a pole of order $3$ there (\cref{thm:poles}).

  \emph{A series for $z / \sin z$.} This is $1/h(z)$, which is analytic near
  $0$. Write $h = 1 - c$ with $c = \frac{z^2}{6} - \frac{z^4}{120} + \dotsb$.
  For small $z$, $\abs{c} < 1$, and the geometric series gives
  \[
    \frac{z}{\sin z} = 1 + c + c^2 + \dotsb
    = 1 + \frac{z^2}{6} + \left(\frac{1}{36} - \frac{1}{120}\right) z^4 + \dotsb
    = 1 + \frac{z^2}{6} + \frac{7 z^4}{360} + \dotsb.
  \]

  \emph{By the formula.} With $m = 3$,
  \[
    \Res(f; 0)
    = \frac{1}{2!} \lim_{z \to 0} \frac{\dd^2}{\dd z^2} \left[ \frac{z}{\sin z} \right]
    = \frac12 \lim_{z \to 0} \left( \frac13 + \frac{7 z^2}{30} + \dotsb \right)
    = \frac16,
  \]
  differentiating term by term (\cref{thm:termwise}).

  \emph{From the Laurent series.} Dividing by $z^3$,
  \[
    f(z) = \frac{1}{z^3} \cdot \frac{z}{\sin z}
    = \frac{1}{z^3} + \frac{1}{6z} + \frac{7z}{360} + \dotsb,
  \]
  so again $\Res(f; 0) = a_{-1} = \frac16$.

  \emph{The integral.} The singularities of $f$ are the zeros $z = k\pi$ of
  $z^2 \sin z$, and only $z = 0$ lies inside $\abs{z} = 1$. By
  \cref{thm:residue-theorem},
  \[
    \oint_{\abs{z} = 1} \frac{\dd z}{z^2 \sin z}
    = 2\pi i \cdot \frac16 = \frac{\pi i}{3}. \qedhere
  \]
\end{solution}

\begin{exercise}\label{ex:residues-two-functions}
  Find the singularities and the residues there of
  \[
    \text{(a)}\ f(z) = \frac{\sin z}{z(z + 1)},
    \qquad
    \text{(b)}\ g(z) = \frac{\cos z}{z (z - 1)^2}.
  \]
\end{exercise}

\begin{solution}
  \emph{(a) At $0$.} Since $\frac{\sin z}{z} = 1 - \frac{z^2}{3!} + \dotsb$,
  $f(z) \to 1$ as $z \to 0$. The singularity is removable
  (\cref{thm:removable}), so there are no negative powers and $\Res(f; 0) = 0$.

  \emph{(a) At $-1$.} $f(z) = \frac{\sin z / z}{z + 1}$, and
  $\frac{\sin(-1)}{-1} = \sin 1 \neq 0$, so $-1$ is a simple pole and
  \[
    \Res(f; -1) = \lim_{z \to -1} \frac{\sin z}{z} = \sin 1.
  \]

  \emph{(b) At $0$.} $g(z) = \frac{\cos z / (z - 1)^2}{z}$, with numerator
  $1 \neq 0$ at $z = 0$. This is a simple pole, and
  \[
    \Res(g; 0) = \lim_{z \to 0} \frac{\cos z}{(z - 1)^2} = 1.
  \]

  \emph{(b) At $1$.} $g(z) = \frac{\cos z / z}{(z - 1)^2}$, with numerator
  $\cos 1 \neq 0$ at $z = 1$, so this is a pole of order $2$. By
  \cref{thm:residue-at-pole} with $m = 2$,
  \[
    \Res(g; 1)
    = \lim_{z \to 1} \frac{\dd}{\dd z} \frac{\cos z}{z}
    = \lim_{z \to 1} \frac{-z \sin z - \cos z}{z^2}
    = -\sin 1 - \cos 1. \qedhere
  \]
\end{solution}
